% Fichier engendré par tools/generate-tex.py à partir de 02-geometrie-plane.
% Ne pas éditer à la main : tout le texte provient des commentaires du fichier Lean.
\documentclass[11pt,a4paper]{article}

% Compile aussi bien avec pdflatex qu'avec xelatex ou tectonic.
\usepackage{iftex}
\ifPDFTeX
  \usepackage[T1]{fontenc}
  \usepackage[utf8]{inputenc}
\else
  \usepackage{fontspec}
\fi
\usepackage[french]{babel}
\usepackage{amsmath,amssymb,amsthm}
\usepackage[margin=2.5cm]{geometry}
\usepackage[hidelinks]{hyperref}

\theoremstyle{definition}
\newtheorem{definition}{Définition}
\theoremstyle{plain}
\newtheorem{theoreme}{Théorème}
\newtheorem{lemme}[theoreme]{Lemme}

% Renvoi à la déclaration Lean, une ligne après la démonstration, aligné à droite.
\newcommand{\source}[2]{\par\smallskip\noindent\hfill{\footnotesize\href{#1}{\texttt{#2}}}\par\smallskip}

\title{Géométrie plane}
\date{}

\begin{document}
\maketitle
% Les identifiants Lean en machine à écrire ne se coupent pas : on tolère des
% espacements plus lâches plutôt que des lignes qui débordent.
\sloppy

\section{AnglesOrientes}

\noindent
Collège \textemdash{} angles alternes-internes et correspondants.

Ces énoncés demandent des angles *orientés* : deux angles alternes-internes ont même
mesure non orientée sans être le même angle, et la réciproque \textemdash{} l'égalité des angles
entraîne le parallélisme \textemdash{} est fausse si l'on oublie l'orientation. On travaille donc dans
un plan orienté, ce qu'impose l'hypothèse \texttt{Fact (finrank \(\mathbb{R}\) V = 2)} : elle vaut pour tout
le fichier, d'où sa séparation d'avec GeometriePlane.lean.
Énoncés et démonstrations en français : voir GeometriePlane.tex.

\subsection{Deux parallèles coupées par une sécante : angles correspondants}

\begin{theoreme}
Deux droites parallèles coupées par une sécante font avec elle des \emph{angles
correspondants} égaux : multiplier le vecteur directeur par un réel strictement positif —
c'est-à-dire prendre une parallèle de même sens — ne change pas l'angle orienté avec la
sécante.
\end{theoreme}

\begin{proof}
C'est la propriété de l'angle orienté : $\widehat{(ku, w)} = \widehat{(u, w)}$ pour
$k > 0$, le vecteur $ku$ pointant dans la même direction que $u$.
\end{proof}
\source{https://github.com/Commutator-IO/learn-lean/blob/main/courses/01-college/02-geometrie-plane/AnglesOrientes.lean\#L23}{AnglesOrientes.lean\#L23}

\subsection{Angles alternes-internes}

\begin{theoreme}
Les \emph{angles alternes-internes} sont égaux : ils se déduisent des angles
correspondants en retournant les deux demi-droites,
\[ \widehat{(ku, w)} = \widehat{(-u, -w)} \qquad (k > 0) . \]
\end{theoreme}

\begin{proof}
On applique l'énoncé précédent, puis le fait qu'un angle orienté ne change pas quand on
remplace ses deux vecteurs par leurs opposés — ce qui revient à faire tourner la figure
d'un demi-tour.
\end{proof}
\source{https://github.com/Commutator-IO/learn-lean/blob/main/courses/01-college/02-geometrie-plane/AnglesOrientes.lean\#L31}{AnglesOrientes.lean\#L31}

\subsection{Réciproque : égalité des angles alternes-internes \(\implies\) parallélisme}

\begin{theoreme}
Réciproque : si deux droites font le même angle orienté avec une sécante, leurs
directions sont colinéaires — elles sont parallèles.
\end{theoreme}

\begin{proof}
La relation de Chasles sur les angles orientés donne
\[ \widehat{(u, v)} = \widehat{(u, w)} + \widehat{(w, v)} , \]
et l'hypothèse remplace le premier terme par $\widehat{(v, w)}$, dont la somme avec
$\widehat{(w, v)}$ est nulle. L'angle $\widehat{(u, v)}$ est donc nul, ce qui caractérise
deux vecteurs colinéaires de même sens.
\end{proof}
\source{https://github.com/Commutator-IO/learn-lean/blob/main/courses/01-college/02-geometrie-plane/AnglesOrientes.lean\#L40}{AnglesOrientes.lean\#L40}

\begin{theoreme}
Sous la forme du cours : l'égalité des angles alternes-internes entraîne le
parallélisme.
\end{theoreme}

\begin{proof}
On se ramène à l'énoncé précédent en retournant les deux demi-droites, ce qui ne change
pas l'angle orienté.
\end{proof}
\source{https://github.com/Commutator-IO/learn-lean/blob/main/courses/01-college/02-geometrie-plane/AnglesOrientes.lean\#L53}{AnglesOrientes.lean\#L53}

\section{GeometriePlane}

\noindent
Collège \textemdash{} section \og{} Géométrie plane \fg{}.
Le plan est un espace vectoriel euclidien : un point est un vecteur, une droite est
l'ensemble des \texttt{a + t • u}. Les angles sont ceux de Mathlib, non orientés, mesurés en
radians : l'angle plat vaut \(\pi\) et l'angle droit \(\pi\)/2.
Énoncés et démonstrations en français : voir GeometriePlane.tex.

\subsection{Définitions}

\noindent
Toutes les définitions du fichier sont posées ici, avant les démonstrations.

\begin{definition}
La \emph{droite} passant par $a$ et dirigée par $u$ est l'ensemble
\[ \mathcal{D}(a, u) = \{\, a + tu \ ;\ t \in \mathbb{R} \,\} . \]
\end{definition}
\source{https://github.com/Commutator-IO/learn-lean/blob/main/courses/01-college/02-geometrie-plane/GeometriePlane.lean\#L21}{GeometriePlane.lean\#L21}

\begin{definition}
La \emph{droite passant par deux points} $a$ et $b$ est $\mathcal{D}(a,\, b - a)$.
\end{definition}
\source{https://github.com/Commutator-IO/learn-lean/blob/main/courses/01-college/02-geometrie-plane/GeometriePlane.lean\#L24}{GeometriePlane.lean\#L24}

\begin{definition}
Deux directions sont \emph{parallèles} lorsqu'elles sont colinéaires : $v = ku$ pour un
réel $k$.
\end{definition}
\source{https://github.com/Commutator-IO/learn-lean/blob/main/courses/01-college/02-geometrie-plane/GeometriePlane.lean\#L27}{GeometriePlane.lean\#L27}

\begin{definition}
Deux directions sont \emph{perpendiculaires} lorsque leur produit scalaire est nul :
$\langle u, v\rangle = 0$.
\end{definition}
\source{https://github.com/Commutator-IO/learn-lean/blob/main/courses/01-college/02-geometrie-plane/GeometriePlane.lean\#L30}{GeometriePlane.lean\#L30}

\begin{definition}
Le \emph{milieu} de $[ab]$ est $\dfrac{a + b}{2}$.
\end{definition}
\source{https://github.com/Commutator-IO/learn-lean/blob/main/courses/01-college/02-geometrie-plane/GeometriePlane.lean\#L33}{GeometriePlane.lean\#L33}

\begin{definition}
Un quadrilatère $abcd$ est un \emph{parallélogramme} lorsque deux côtés opposés portent
le même vecteur : $b - a = c - d$.
\end{definition}
\source{https://github.com/Commutator-IO/learn-lean/blob/main/courses/01-college/02-geometrie-plane/GeometriePlane.lean\#L37}{GeometriePlane.lean\#L37}

\begin{definition}
La \emph{distance d'un point à une droite} passant par $a$ et dirigée par un vecteur
unitaire $u$ est la longueur de la composante orthogonale :
\[ d(p, \mathcal{D}) = \bigl\| (p - a) - \langle p - a, u\rangle\, u \bigr\| . \]
\end{definition}
\source{https://github.com/Commutator-IO/learn-lean/blob/main/courses/01-college/02-geometrie-plane/GeometriePlane.lean\#L41}{GeometriePlane.lean\#L41}

\begin{definition}
Le \emph{centre de gravité} du triangle $abc$ est $\dfrac{a + b + c}{3}$.
\end{definition}
\source{https://github.com/Commutator-IO/learn-lean/blob/main/courses/01-college/02-geometrie-plane/GeometriePlane.lean\#L45}{GeometriePlane.lean\#L45}

\subsection{Par deux points distincts passe une droite et une seule}

\begin{theoreme}
Les deux points appartiennent à la droite qu'ils définissent.
\end{theoreme}

\begin{proof}
Le paramètre $t = 0$ donne $a$, le paramètre $t = 1$ donne $a + (b - a) = b$.
\end{proof}
\source{https://github.com/Commutator-IO/learn-lean/blob/main/courses/01-college/02-geometrie-plane/GeometriePlane.lean\#L50}{GeometriePlane.lean\#L50}

\begin{theoreme}
Et c'est la seule : toute droite contenant $a$ et $b$ distincts contient la droite
$(ab)$.
\end{theoreme}

\begin{proof}
Si $a = c + su$ et $b = c + tu$, alors tout point $a + r(b - a)$ de la droite $(ab)$
s'écrit
\[ c + su + r\bigl((c + tu) - (c + su)\bigr) = c + \bigl(s + r(t - s)\bigr)u , \]
qui est bien un point de la droite $\mathcal{D}(c, u)$.
\end{proof}
\source{https://github.com/Commutator-IO/learn-lean/blob/main/courses/01-college/02-geometrie-plane/GeometriePlane.lean\#L54}{GeometriePlane.lean\#L54}

\subsection{Perpendicularité et parallélisme}

\begin{theoreme}
Si deux droites sont parallèles, toute perpendiculaire à l'une est perpendiculaire à
l'autre.
\end{theoreme}

\begin{proof}
Si $v = ku$ et $\langle w, u\rangle = 0$, alors
$\langle w, v\rangle = k\,\langle w, u\rangle = 0$.
\end{proof}
\source{https://github.com/Commutator-IO/learn-lean/blob/main/courses/01-college/02-geometrie-plane/GeometriePlane.lean\#L67}{GeometriePlane.lean\#L67}

\begin{theoreme}
Deux droites parallèles à une même droite sont parallèles entre elles.
\end{theoreme}

\begin{proof}
Si $v = ku$ et $w = \ell v$, alors $w = (\ell k)u$ : la colinéarité se transmet, les
coefficients se multipliant.
\end{proof}
\source{https://github.com/Commutator-IO/learn-lean/blob/main/courses/01-college/02-geometrie-plane/GeometriePlane.lean\#L74}{GeometriePlane.lean\#L74}

\begin{theoreme}
Dans le plan, deux droites perpendiculaires à une même droite sont parallèles entre
elles : leur déterminant est nul,
\[ u_0 v_1 - u_1 v_0 = 0 . \]
\end{theoreme}

\begin{proof}
Les deux conditions d'orthogonalité s'écrivent en coordonnées
\[ w_0 u_0 + w_1 u_1 = 0, \qquad w_0 v_0 + w_1 v_1 = 0 , \]
et $w$ n'étant pas nul, l'une de ses coordonnées ne l'est pas. Supposons $w_0 \ne 0$ : on
en tire $u_0 = -w_1 u_1 / w_0$ et $v_0 = -w_1 v_1 / w_0$, d'où
\[ u_0 v_1 - u_1 v_0 = \frac{-w_1 u_1 v_1 + w_1 u_1 v_1}{w_0} = 0 . \]
Le cas $w_1 \ne 0$ est symétrique.

La dimension est essentielle : dans l'espace, deux perpendiculaires à une même droite ne
sont pas parallèles en général. C'est pourquoi cet énoncé, seul du fichier, est posé dans
le plan de coordonnées et non dans un espace euclidien quelconque.
\end{proof}
\source{https://github.com/Commutator-IO/learn-lean/blob/main/courses/01-college/02-geometrie-plane/GeometriePlane.lean\#L83}{GeometriePlane.lean\#L83}

\begin{theoreme}
Dans le plan, un déterminant nul signifie que les deux vecteurs sont colinéaires : si
$u \ne 0$ et $u_0 v_1 - u_1 v_0 = 0$, alors $v = ku$ pour un réel $k$.
\end{theoreme}

\begin{proof}
L'une des coordonnées de $u$ est non nulle ; supposons $u_0 \ne 0$ et posons
$k = v_0/u_0$. La première coordonnée de $ku$ vaut alors $v_0$ par construction, et la
seconde
\[ \frac{v_0 u_1}{u_0} = \frac{u_0 v_1}{u_0} = v_1 \]
grâce à la nullité du déterminant. Le cas $u_1 \ne 0$ est symétrique.
\end{proof}
\source{https://github.com/Commutator-IO/learn-lean/blob/main/courses/01-college/02-geometrie-plane/GeometriePlane.lean\#L110}{GeometriePlane.lean\#L110}

\begin{theoreme}
Sous la forme du cours : deux droites perpendiculaires à une même droite sont
parallèles entre elles.
\end{theoreme}

\begin{proof}
Leurs directions ont un déterminant nul d'après l'énoncé précédent : elles sont donc
colinéaires, et les droites parallèles.
\end{proof}
\source{https://github.com/Commutator-IO/learn-lean/blob/main/courses/01-college/02-geometrie-plane/GeometriePlane.lean\#L128}{GeometriePlane.lean\#L128}

\subsection{Le plus court chemin d'un point à une droite est le segment perpendiculaire}

\begin{theoreme}
Le plus court chemin d'un point à une droite est le segment perpendiculaire : si
$p - h$ est perpendiculaire à la direction de la droite, alors
\[ \|p - h\| \le \|p - (h + tu)\| \quad \text{pour tout } t . \]
\end{theoreme}

\begin{proof}
On développe le carré de la distance à un point courant de la droite :
\[ \|p - (h + tu)\|^2 = \|(p - h) - tu\|^2
   = \|p - h\|^2 - 2t\,\langle p - h, u\rangle + t^2\|u\|^2 . \]
Le terme du milieu est nul par hypothèse, et le dernier est positif : la distance est
minimale pour $t = 0$, c'est-à-dire au pied de la perpendiculaire.
\end{proof}
\source{https://github.com/Commutator-IO/learn-lean/blob/main/courses/01-college/02-geometrie-plane/GeometriePlane.lean\#L137}{GeometriePlane.lean\#L137}

\subsection{Angles opposés par le sommet, adjacents, supplémentaires}

\begin{theoreme}
Deux angles opposés par le sommet sont égaux.
\end{theoreme}

\begin{proof}
Les deux couples de vecteurs sont opposés, et l'angle de deux vecteurs ne change pas
quand on les remplace tous deux par leurs opposés.
\end{proof}
\source{https://github.com/Commutator-IO/learn-lean/blob/main/courses/01-college/02-geometrie-plane/GeometriePlane.lean\#L150}{GeometriePlane.lean\#L150}

\begin{theoreme}
Deux angles adjacents portés par des demi-droites opposées sont supplémentaires :
\[ \widehat{(u, v)} + \widehat{(u, -v)} = \pi . \]
\end{theoreme}

\begin{proof}
Remplacer un vecteur par son opposé remplace l'angle par son supplémentaire :
$\widehat{(u, -v)} = \pi - \widehat{(u, v)}$. La somme vaut donc $\pi$, l'angle plat.
\end{proof}
\source{https://github.com/Commutator-IO/learn-lean/blob/main/courses/01-college/02-geometrie-plane/GeometriePlane.lean\#L156}{GeometriePlane.lean\#L156}

\subsection{Caractérisation de la médiatrice}

\begin{theoreme}
Caractérisation de la médiatrice : un point est équidistant de $a$ et $b$ si et
seulement s'il appartient à la médiatrice de $[ab]$.
\end{theoreme}

\begin{proof}
C'est la définition de la médiatrice dans la bibliothèque, à l'ordre des arguments de
la distance près.
\end{proof}
\source{https://github.com/Commutator-IO/learn-lean/blob/main/courses/01-college/02-geometrie-plane/GeometriePlane.lean\#L165}{GeometriePlane.lean\#L165}

\subsection{Somme des angles d'un triangle}

\begin{theoreme}
La somme des angles d'un triangle vaut l'angle plat :
\[ \widehat{abc} + \widehat{bca} + \widehat{cab} = \pi . \]
\end{theoreme}

\begin{proof}
C'est le théorème de la bibliothèque, valable dès que deux des sommets sont distincts.
Sa démonstration passe par la trigonométrie du produit scalaire, là où le collège trace
une parallèle et utilise les angles alternes-internes.
\end{proof}
\source{https://github.com/Commutator-IO/learn-lean/blob/main/courses/01-college/02-geometrie-plane/GeometriePlane.lean\#L173}{GeometriePlane.lean\#L173}

\subsection{Triangle isocèle : angles à la base}

\begin{theoreme}
Dans un triangle isocèle, les angles à la base sont égaux :
$\mathrm{dist}(a,b) = \mathrm{dist}(a,c) \Rightarrow \widehat{abc} = \widehat{acb}$.
\end{theoreme}

\begin{proof}
Deux côtés égaux rendent la figure symétrique par rapport à la médiatrice du troisième,
qui échange les deux angles à la base.
\end{proof}
\source{https://github.com/Commutator-IO/learn-lean/blob/main/courses/01-college/02-geometrie-plane/GeometriePlane.lean\#L180}{GeometriePlane.lean\#L180}

\begin{theoreme}
Réciproque : si les angles à la base sont égaux, le triangle est isocèle.
\end{theoreme}

\begin{proof}
C'est le résultat réciproque de la bibliothèque, valable dès que le triangle n'est pas
aplati — d'où l'hypothèse que l'angle au sommet n'est pas l'angle plat.
\end{proof}
\source{https://github.com/Commutator-IO/learn-lean/blob/main/courses/01-college/02-geometrie-plane/GeometriePlane.lean\#L184}{GeometriePlane.lean\#L184}

\begin{theoreme}
Dans un triangle équilatéral, les trois angles valent $60$ degrés, soit $\pi/3$.
\end{theoreme}

\begin{proof}
Les deux égalités de côtés donnent, par le théorème de l'isocèle appliqué deux fois,
\[ \widehat{abc} = \widehat{acb}, \qquad \widehat{bac} = \widehat{bca} . \]
Un angle ne changeant pas quand on échange ses deux extrémités, les trois angles du
triangle sont donc égaux ; leur somme valant $\pi$, chacun vaut $\pi/3$.
\end{proof}
\source{https://github.com/Commutator-IO/learn-lean/blob/main/courses/01-college/02-geometrie-plane/GeometriePlane.lean\#L189}{GeometriePlane.lean\#L189}

\subsection{Inégalité triangulaire}

\begin{theoreme}
Inégalité triangulaire : $\mathrm{dist}(a,c) \le \mathrm{dist}(a,b) + \mathrm{dist}(b,c)$.
\end{theoreme}

\begin{proof}
C'est l'inégalité triangulaire de la distance : le chemin direct est le plus court.
\end{proof}
\source{https://github.com/Commutator-IO/learn-lean/blob/main/courses/01-college/02-geometrie-plane/GeometriePlane.lean\#L204}{GeometriePlane.lean\#L204}

\begin{theoreme}
Le cas d'égalité caractérise l'alignement :
\[ \mathrm{dist}(a,b) + \mathrm{dist}(b,c) = \mathrm{dist}(a,c)
   \iff b \text{ appartient au segment } [ac] . \]
\end{theoreme}

\begin{proof}
L'égalité n'a lieu que si le détour par $b$ ne coûte rien, c'est-à-dire si $b$ est entre
$a$ et $c$. Le résultat vaut dans un espace strictement convexe, ce qu'est tout espace
euclidien.
\end{proof}
\source{https://github.com/Commutator-IO/learn-lean/blob/main/courses/01-college/02-geometrie-plane/GeometriePlane.lean\#L208}{GeometriePlane.lean\#L208}

\subsection{Théorème de Pythagore, sa réciproque et sa contraposée}

\begin{theoreme}
Théorème de Pythagore et sa réciproque :
\[ \mathrm{dist}(a,c)^2 = \mathrm{dist}(a,b)^2 + \mathrm{dist}(c,b)^2
   \iff \widehat{abc} = \frac{\pi}{2} . \]
\end{theoreme}

\begin{proof}
Le carré de la distance se développe en produits scalaires :
\[ \|a - c\|^2 = \|a - b\|^2 + \|c - b\|^2 - 2\,\langle a - b, c - b\rangle , \]
et le double produit est nul exactement quand les deux côtés sont perpendiculaires,
c'est-à-dire quand l'angle en $b$ est droit. L'équivalence donne d'un coup le théorème et
sa réciproque.
\end{proof}
\source{https://github.com/Commutator-IO/learn-lean/blob/main/courses/01-college/02-geometrie-plane/GeometriePlane.lean\#L215}{GeometriePlane.lean\#L215}

\begin{theoreme}
Contraposée : si l'égalité des carrés est en défaut, le triangle n'est pas rectangle.
\end{theoreme}

\begin{proof}
C'est la contraposée de la réciproque : si l'angle était droit, l'égalité aurait lieu.
Cette forme est celle qu'on utilise en classe pour montrer qu'un triangle \emph{n'est
pas} rectangle.
\end{proof}
\source{https://github.com/Commutator-IO/learn-lean/blob/main/courses/01-college/02-geometrie-plane/GeometriePlane.lean\#L220}{GeometriePlane.lean\#L220}

\subsection{Théorème des milieux et sa réciproque}

\begin{theoreme}
Théorème des milieux : la droite des milieux est parallèle au troisième côté,
\[ \text{milieu}(a,c) - \text{milieu}(a,b) = \tfrac{1}{2}(c - b) . \]
\end{theoreme}

\begin{proof}
On calcule directement :
\[ \frac{a+c}{2} - \frac{a+b}{2} = \frac{c - b}{2} . \]
Le vecteur obtenu est colinéaire à $c - b$ : les deux droites sont parallèles.
\end{proof}
\source{https://github.com/Commutator-IO/learn-lean/blob/main/courses/01-college/02-geometrie-plane/GeometriePlane.lean\#L227}{GeometriePlane.lean\#L227}

\begin{theoreme}
Conséquence sur les longueurs : le segment des milieux mesure la moitié du troisième
côté.
\end{theoreme}

\begin{proof}
La norme du vecteur précédent vaut $\tfrac12\|b - c\|$, le facteur $\tfrac12$ sortant de
la norme en valeur absolue.
\end{proof}
\source{https://github.com/Commutator-IO/learn-lean/blob/main/courses/01-college/02-geometrie-plane/GeometriePlane.lean\#L234}{GeometriePlane.lean\#L234}

\begin{theoreme}
Réciproque : le point du côté $[ac]$ dont l'écart au milieu de $[ab]$ vaut
$\tfrac12(c-b)$ est le milieu de $[ac]$.
\end{theoreme}

\begin{proof}
En reportant l'écart depuis le milieu de $[ab]$ :
\[ \frac{a+b}{2} + \frac{c-b}{2} = \frac{a+c}{2} , \]
qui est le milieu de $[ac]$.
\end{proof}
\source{https://github.com/Commutator-IO/learn-lean/blob/main/courses/01-college/02-geometrie-plane/GeometriePlane.lean\#L246}{GeometriePlane.lean\#L246}

\subsection{Théorème de Thalès et sa réciproque}

\begin{theoreme}
Théorème de Thalès : si $m$ et $n$ divisent $[ab]$ et $[ac]$ dans le même rapport $t$,
alors
\[ \bigl(a + t(b-a)\bigr) - \bigl(a + t(c-a)\bigr) = t\,(b - c) . \]
\end{theoreme}

\begin{proof}
Le calcul est immédiat : les deux $a$ se simplifient et le facteur $t$ se met en
facteur. Le vecteur $mn$ est donc colinéaire à $bc$ — les droites sont parallèles — et le
coefficient est le rapport commun.
\end{proof}
\source{https://github.com/Commutator-IO/learn-lean/blob/main/courses/01-college/02-geometrie-plane/GeometriePlane.lean\#L257}{GeometriePlane.lean\#L257}

\begin{theoreme}
Conséquence sur les longueurs : les trois rapports sont égaux,
\[ \mathrm{dist}(m, n) = t \times \mathrm{dist}(b, c) . \]
\end{theoreme}

\begin{proof}
On prend la norme de l'égalité précédente, le facteur $t \ge 0$ sortant tel quel.
\end{proof}
\source{https://github.com/Commutator-IO/learn-lean/blob/main/courses/01-college/02-geometrie-plane/GeometriePlane.lean\#L262}{GeometriePlane.lean\#L262}

\begin{theoreme}
Réciproque de Thalès : si $mn$ est parallèle à $bc$ avec le rapport $t$, alors $n$
divise $[ac]$ dans ce même rapport.
\end{theoreme}

\begin{proof}
En reportant l'écart depuis $m = a + t(b-a)$ :
\[ n = a + t(b-a) + t(c-b) = a + t(c-a) . \]
Le point $n$ divise donc $[ac]$ dans le rapport $t$.
\end{proof}
\source{https://github.com/Commutator-IO/learn-lean/blob/main/courses/01-college/02-geometrie-plane/GeometriePlane.lean\#L268}{GeometriePlane.lean\#L268}

\subsection{Relations trigonométriques}

\begin{theoreme}
Relation fondamentale de la trigonométrie : $\cos^2 x + \sin^2 x = 1$.
\end{theoreme}

\begin{proof}
C'est l'identité de Pythagore sur le cercle trigonométrique : le point de coordonnées
$(\cos x, \sin x)$ est à distance $1$ de l'origine.
\end{proof}
\source{https://github.com/Commutator-IO/learn-lean/blob/main/courses/01-college/02-geometrie-plane/GeometriePlane.lean\#L278}{GeometriePlane.lean\#L278}

\begin{theoreme}
La tangente est le quotient du sinus par le cosinus : $\tan x = \dfrac{\sin x}{\cos x}$.
\end{theoreme}

\begin{proof}
C'est la définition de la tangente.
\end{proof}
\source{https://github.com/Commutator-IO/learn-lean/blob/main/courses/01-college/02-geometrie-plane/GeometriePlane.lean\#L282}{GeometriePlane.lean\#L282}

\subsection{Cercle circonscrit}

\begin{theoreme}
Les médiatrices d'un triangle sont concourantes : il existe un point équidistant des
trois sommets, centre du cercle circonscrit.
\end{theoreme}

\begin{proof}
On prend le centre du cercle circonscrit au triangle, dont la bibliothèque garantit
l'existence dès que les trois points ne sont pas alignés. Il est à égale distance des trois
sommets — cette distance étant le rayon —, donc situé sur chacune des trois médiatrices,
qui sont ainsi concourantes.
\end{proof}
\source{https://github.com/Commutator-IO/learn-lean/blob/main/courses/01-college/02-geometrie-plane/GeometriePlane.lean\#L288}{GeometriePlane.lean\#L288}

\subsection{Triangle inscrit dans un demi-cercle}

\begin{theoreme}
Un triangle inscrit dans un cercle dont un côté est un diamètre est rectangle, et
réciproquement :
\[ \widehat{abc} = \frac{\pi}{2} \iff b \text{ appartient au cercle de diamètre } [ac] . \]
\end{theoreme}

\begin{proof}
C'est le théorème de Thalès sur le cercle, tel que la bibliothèque l'établit. L'énoncé
donne les deux sens à la fois : le triangle inscrit dans un demi-cercle est rectangle, et
l'hypoténuse d'un triangle rectangle est un diamètre de son cercle circonscrit.
\end{proof}
\source{https://github.com/Commutator-IO/learn-lean/blob/main/courses/01-college/02-geometrie-plane/GeometriePlane.lean\#L301}{GeometriePlane.lean\#L301}

\begin{theoreme}
Réciproque : l'hypoténuse d'un triangle rectangle est un diamètre du cercle qui passe
par les trois sommets.
\end{theoreme}

\begin{proof}
C'est l'autre sens de l'équivalence de Thalès sur le cercle : l'angle en $b$ étant
droit, $b$ appartient au cercle de diamètre $[ac]$.
\end{proof}
\source{https://github.com/Commutator-IO/learn-lean/blob/main/courses/01-college/02-geometrie-plane/GeometriePlane.lean\#L307}{GeometriePlane.lean\#L307}

\begin{theoreme}
Dans un triangle rectangle, la médiane issue de l'angle droit vaut la moitié de
l'hypoténuse :
\[ \mathrm{dist}\bigl(b, \text{milieu}(a,c)\bigr) = \frac{\mathrm{dist}(a,c)}{2} . \]
\end{theoreme}

\begin{proof}
L'angle droit en $b$ donne $\langle a - b, c - b\rangle = 0$. Or
\[ b - \text{milieu}(a,c) = \tfrac{1}{2}\bigl((b-a) + (b-c)\bigr) , \]
et l'orthogonalité donne, d'une part
\[ \|(b-a) + (b-c)\|^2 = \|b-a\|^2 + \|b-c\|^2 , \]
d'autre part, en écrivant $a - c = (b-c) - (b-a)$,
\[ \|a-c\|^2 = \|b-a\|^2 + \|b-c\|^2 . \]
Les deux normes sont donc égales, et la médiane vaut la moitié de $\|a - c\|$.
\end{proof}
\source{https://github.com/Commutator-IO/learn-lean/blob/main/courses/01-college/02-geometrie-plane/GeometriePlane.lean\#L313}{GeometriePlane.lean\#L313}

\subsection{Tangente à un cercle}

\begin{theoreme}
La droite passant par un point du cercle et perpendiculaire au rayon en ce point ne
rencontre le cercle qu'en ce point : c'est la tangente.
\end{theoreme}

\begin{proof}
Si $p + tu$ est à la même distance du centre que $p$, alors
\[ \|p - o\|^2 = \|(p-o) + tu\|^2 = \|p-o\|^2 + t^2\|u\|^2 , \]
le double produit étant nul par perpendicularité. Il reste $t^2\|u\|^2 = 0$, et comme
$u \ne 0$, on obtient $t = 0$ : le seul point commun est $p$.
\end{proof}
\source{https://github.com/Commutator-IO/learn-lean/blob/main/courses/01-college/02-geometrie-plane/GeometriePlane.lean\#L341}{GeometriePlane.lean\#L341}

\subsection{Parallélogramme}

\begin{theoreme}
Un quadrilatère est un parallélogramme si et seulement si ses diagonales se coupent en
leur milieu : $b - a = c - d \iff \text{milieu}(a,c) = \text{milieu}(b,d)$.
\end{theoreme}

\begin{proof}
Les deux conditions sont deux écritures de la même égalité $a + c = b + d$ : la
première la donne en retranchant, la seconde en divisant par deux.
\end{proof}
\source{https://github.com/Commutator-IO/learn-lean/blob/main/courses/01-college/02-geometrie-plane/GeometriePlane.lean\#L364}{GeometriePlane.lean\#L364}

\begin{theoreme}
Dans un parallélogramme, les côtés opposés ont la même longueur.
\end{theoreme}

\begin{proof}
De $b - a = c - d$ on tire $a - b = d - c$ : les deux côtés portent des vecteurs opposés,
donc de même norme.
\end{proof}
\source{https://github.com/Commutator-IO/learn-lean/blob/main/courses/01-college/02-geometrie-plane/GeometriePlane.lean\#L376}{GeometriePlane.lean\#L376}

\begin{theoreme}
Réciproque, sous l'hypothèse que le quadrilatère n'est pas croisé : si les côtés
opposés ont même longueur et pointent dans le même sens, le quadrilatère est un
parallélogramme.
\end{theoreme}

\begin{proof}
Deux vecteurs de même sens et de même norme sont égaux. L'hypothèse de même sens est
celle que la figure porte au collège et que l'énoncé tait ; l'énoncé suivant montre qu'elle
ne peut pas être omise.
\end{proof}
\source{https://github.com/Commutator-IO/learn-lean/blob/main/courses/01-college/02-geometrie-plane/GeometriePlane.lean\#L390}{GeometriePlane.lean\#L390}

\begin{theoreme}
Sans cette hypothèse, l'énoncé est faux : le quadrilatère $(0, u, 0, u)$ a ses côtés
opposés de même longueur sans être un parallélogramme.
\end{theoreme}

\begin{proof}
Les deux côtés opposés sont $u$ et $-u$ : de même longueur, mais de sens contraires.
L'égalité $b - a = c - d$ donnerait $2u = 0$, donc $u = 0$, exclu. C'est le quadrilatère
croisé, que le dessin écarte sans que le texte le dise.
\end{proof}
\source{https://github.com/Commutator-IO/learn-lean/blob/main/courses/01-college/02-geometrie-plane/GeometriePlane.lean\#L401}{GeometriePlane.lean\#L401}

\begin{theoreme}
Dans un parallélogramme, les angles opposés sont égaux.
\end{theoreme}

\begin{proof}
L'égalité $b - a = c - d$ équivaut à $b + d = c + a$, d'où
\[ d - c = -(b - a), \qquad b - c = -(d - a) . \]
Les deux couples de vecteurs portant les angles en $a$ et en $c$ sont donc opposés, et
l'angle est inchangé quand on remplace les deux vecteurs par leurs opposés.
\end{proof}
\source{https://github.com/Commutator-IO/learn-lean/blob/main/courses/01-college/02-geometrie-plane/GeometriePlane.lean\#L412}{GeometriePlane.lean\#L412}

\subsection{Concours des médianes}

\begin{theoreme}
Le centre de gravité est sur la médiane issue de $a$.
\end{theoreme}

\begin{proof}
La médiane issue de $a$ joint $a$ au milieu de $[bc]$. Le paramètre $t = \tfrac23$ y
donne
\[ a + \frac{2}{3}\left(\frac{b+c}{2} - a\right) = \frac{a+b+c}{3} , \]
qui est le centre de gravité : c'est la propriété du tiers, le centre de gravité étant aux
deux tiers de la médiane en partant du sommet.
\end{proof}
\source{https://github.com/Commutator-IO/learn-lean/blob/main/courses/01-college/02-geometrie-plane/GeometriePlane.lean\#L428}{GeometriePlane.lean\#L428}

\begin{theoreme}
Sur celle issue de $b$.
\end{theoreme}

\begin{proof}
Même calcul, l'expression du centre de gravité étant symétrique en les trois sommets.
\end{proof}
\source{https://github.com/Commutator-IO/learn-lean/blob/main/courses/01-college/02-geometrie-plane/GeometriePlane.lean\#L433}{GeometriePlane.lean\#L433}

\begin{theoreme}
Et sur celle issue de $c$ : les trois médianes sont donc concourantes en ce point.
\end{theoreme}

\begin{proof}
Même calcul une troisième fois : les trois médianes passent par le centre de gravité.
\end{proof}
\source{https://github.com/Commutator-IO/learn-lean/blob/main/courses/01-college/02-geometrie-plane/GeometriePlane.lean\#L438}{GeometriePlane.lean\#L438}

\subsection{Caractérisation de la bissectrice : équidistance aux deux côtés}

\begin{theoreme}
Caractérisation de la bissectrice : les points de la bissectrice d'un angle sont à
égale distance de ses deux côtés. La bissectrice est portée par la somme des deux vecteurs
unitaires — traduction vectorielle de « elle partage l'angle en deux ».
\end{theoreme}

\begin{proof}
Soit $p = a + t(u + v)$ un point de la bissectrice, $u$ et $v$ étant unitaires. La
composante de $p - a$ orthogonale au premier côté vaut
\[ t(u+v) - t\bigl(1 + \langle v, u\rangle\bigr)u = t\bigl(v - \langle v, u\rangle u\bigr) , \]
et symétriquement pour le second côté. Or ces deux vecteurs ont même norme, puisque
\[ \|v - \langle u,v\rangle u\|^2 = 1 - \langle u,v\rangle^2
   = \|u - \langle u,v\rangle v\|^2 , \]
les deux vecteurs étant unitaires. Les deux distances valent donc toutes deux
$|t|\sqrt{1 - \langle u,v\rangle^2}$.
\end{proof}
\source{https://github.com/Commutator-IO/learn-lean/blob/main/courses/01-college/02-geometrie-plane/GeometriePlane.lean\#L447}{GeometriePlane.lean\#L447}

\subsection{Concours des hauteurs et des bissectrices}

\begin{theoreme}
Les hauteurs d'un triangle sont concourantes : l'orthocentre appartient à chacune
d'elles.
\end{theoreme}

\begin{proof}
C'est le résultat de la bibliothèque : l'orthocentre, défini comme point de Monge du
triangle, appartient à chacune des trois hauteurs. Les trois hauteurs se coupent donc en
un même point.
\end{proof}
\source{https://github.com/Commutator-IO/learn-lean/blob/main/courses/01-college/02-geometrie-plane/GeometriePlane.lean\#L488}{GeometriePlane.lean\#L488}

\begin{theoreme}
Les bissectrices sont concourantes : le centre du cercle inscrit est à la même distance
des trois côtés, cette distance commune étant le rayon du cercle inscrit.
\end{theoreme}

\begin{proof}
La bibliothèque construit le centre du cercle inscrit et établit que sa distance à
chacun des trois côtés vaut le rayon de ce cercle. Être à égale distance de deux côtés,
c'est être sur leur bissectrice : ce point est donc sur les trois bissectrices, qui sont
ainsi concourantes.
\end{proof}
\source{https://github.com/Commutator-IO/learn-lean/blob/main/courses/01-college/02-geometrie-plane/GeometriePlane.lean\#L493}{GeometriePlane.lean\#L493}

\subsection{Repérage : milieu et distance}

\begin{theoreme}
Les coordonnées du milieu sont les moyennes des coordonnées :
$\text{milieu}(a,b)_i = \dfrac{a_i + b_i}{2}$.
\end{theoreme}

\begin{proof}
La définition du milieu s'applique coordonnée par coordonnée.
\end{proof}
\source{https://github.com/Commutator-IO/learn-lean/blob/main/courses/01-college/02-geometrie-plane/GeometriePlane.lean\#L499}{GeometriePlane.lean\#L499}

\begin{theoreme}
La distance entre deux points repérés se calcule par le théorème de Pythagore :
\[ \mathrm{dist}(a,b) = \sqrt{(a_0 - b_0)^2 + (a_1 - b_1)^2} . \]
\end{theoreme}

\begin{proof}
C'est la définition de la distance euclidienne en coordonnées, qui n'est autre que le
théorème de Pythagore appliqué au triangle rectangle dont les côtés sont les écarts en
abscisse et en ordonnée.
\end{proof}
\source{https://github.com/Commutator-IO/learn-lean/blob/main/courses/01-college/02-geometrie-plane/GeometriePlane.lean\#L506}{GeometriePlane.lean\#L506}

\vfill
\noindent\rule{0.35\linewidth}{0.4pt}\par
\noindent{\footnotesize Leonardo de Moura et Sebastian Ullrich,
\emph{The Lean~4 Theorem Prover and Programming Language},
\emph{Automated Deduction — CADE~28}, Springer, 2021, p.~625--635,
\href{https://doi.org/10.1007/978-3-030-79876-5_37}{doi:10.1007/978-3-030-79876-5\_37}.\par}

\end{document}
